\section{Provider-arkitektur}
\label{sec:provider-pattern}

Provider-arkitekturen er det sentrale designvalget i VideoAPI. Behovet oppstod
fordi løsningen skulle tilby ett felles API mot klientene, samtidig som kall
måtte kunne rutes til ulike videomotorer. AntMedia Server og NHN Video har ulike
API-er og domenemodeller: AntMedia eksponerer REST-endepunkter for blant annet
broadcast-ressurser, mens NHN Booking API er orientert rundt opprettelse og
administrasjon av videomøter~\cite{antmedia_rest,nhn_booking}. Samtidig stilte
kravspesifikasjonen krav om at nye videomotorer bør kunne legges til uten at hele
mellomvaren må skrives om~\cite{hdo_kravspec}.

Et alternativ var å samle all backend-logikk i én felles serviceklasse og bruke
betingelser for å skille mellom videomotorene. Dette ville vært enkelt i starten,
men ville gitt tett kobling til de konkrete backendene og gjort løsningen
vanskeligere å utvide. Et annet alternativ var å eksponere separate API-er for
hver backend, for eksempel egne endepunkter for AntMedia og NHN. Dette ville
gjort backend-forskjellene tydelige, men også flyttet kompleksiteten til
klienten, som da måtte vite hvilken videomotor som brukes og hvordan responsene
skal tolkes.

Den valgte løsningen ble derfor å definere et felles
\texttt{IVideoProvider}-grensesnitt. Hver videomotor implementeres som en egen
provider, mens controllerne kun forholder seg til grensesnittet. Designet har
likhetstrekk med adapter-mønsteret, der et felles grensesnitt brukes til å skjule
forskjeller mellom underliggende implementasjoner~\cite{refactoringguru_adapter}.
Riktig provider velges av en egen factory-komponent, slik at controllerne ikke
trenger å kjenne til de konkrete provider-klassene.

Provider-arkitekturen betyr likevel ikke at alle backendene behandles som
identiske. Enkelte operasjoner har ulik betydning i de to systemene. Start- og
stoppoperasjoner passer naturlig med AntMedia sin broadcast-modell, mens NHN
Booking API i større grad beskriver opprettelse og administrasjon av et
videomøte. Provider-laget skal derfor isolere forskjellene, ikke skjule dem helt.
Nye videomotorer kan dermed legges til ved å implementere
\texttt{IVideoProvider}, mens controllerne og klientens API-kontrakt i hovedsak
kan beholdes uendret.

\begin{lstlisting}[language=CSharp, caption={Felles provider-kontrakt}, label={lst:ivideoprovider}]
public interface IVideoProvider
{
    Task<VideoRoom> CreateRoomAsync(CreateRoomRequest request, CancellationToken ct = default);
    Task<VideoRoom> GetRoomAsync(string roomId, CancellationToken ct = default);
    Task DeleteRoomAsync(string roomId, CancellationToken ct = default);
    Task StartStreamAsync(string roomId, CancellationToken ct = default);
    Task StopStreamAsync(string roomId, CancellationToken ct = default);
    Task SetRoomTokenAsync(string roomId, string token, CancellationToken ct = default);
    Task<WebRtcLinkResponse> GenerateWebRtcTokenAsync(string roomId, string type, CancellationToken ct = default);
    Task<bool> IsAvailableAsync(CancellationToken ct = default);
}
\end{lstlisting}
Listing~\ref{lst:ivideoprovider} viser grensesnittet alle videomotorer må
implementere. Dette gjør at controller-laget kan forholde seg til én kontrakt,
mens backend-spesifikk logikk holdes i providerne.